\section{Risiko-analyse}\label{sec:risiko-analyse}

Dette afsnit har til formål at identificere og analysere de væsentligste risici forbundet med at
udvikle Kajakpolo Danmarks managementplatform.
Risici vurderes ud fra sandsynlighed og potentielle konsekvens for projektets tidsplan og funktionalitet.
Punkterne i dette afsnit er informeret af OWASP:2025~(\cite{owasptop10:2025}), overvejelser omkring udviklingsproces baseret på gruppens viden og erfaringer med samme, samt den europæiske GDPR-lovgivning (\cite{GDPR}).

Da udviklingen af systemet gennemføres ved anvendelse af en agil udviklingsmetode, er risikovurderingen også udviklet iterativt
ud fra egen analyse og løbende feedback fra interessenter.

\subsection{Risikoområder}\label{subsec:risikoomrader}
\subsubsection{Uklare krav}
En central risiko er uklare eller ændrede krav, da interessenterne strækker sig over både nationalt- klub- og enkeltbruger niveau.
Projektet arbejder med to hoved-interessenter i form af en kontakt fra Kajakpolo Danmarks bestyrelse samt projektgruppens egen Oliver Rosenkilde,
der er menigt medlem af Aarhus' Kajakpolo-klub.
De skal vurdere applikationen på vegne af den fulde bredde af interessenter, hvilket kan medføre blinde vinkler når de skal repræsentere så bredt.


\subsubsection{Teknisk gæld}
Med projektets agile tilgang er der risiko for at introducere teknisk gæld, hvis hurtige løsninger prioriteres for at nå sprint-mål.
Over tid kan det påvirke systemets vedligeholdbarhed og funktion.
Vi håndterer det ved at afsætte tid i sprintplanlægningen til at refaktorere, og gennem vores proces med at gennemse hinandens pull requests, hvor vedligeholdbarhed og konformitet til designprincip er i fokus.

\subsubsection{Projektorganisation}
Ethvert projekt har organisatoriske risici. I vores tilfælde, hvor gruppen består af kun 3 medlemmer, kan sygdom eller andre forpligtelser for bare ét medlem have en betydelig indvirkning på projektets fremdrift.
Det kan føre til forsinkelser i udviklingen og øget arbejdsbyrde for de resterende medlemmer, hvilket kan påvirke både tidsplanen og kvaliteten af det endelige produkt.
Der bør tages højde for dette ved at have en fleksibel tidsplan og en klar kommunikationsstrategi, så alle medlemmer er opmærksomme på projektets status og kan tilpasse sig ændringer i arbejdsbyrden.
Desuden bør udviklingsopgaverne fordeles jævnt, og der bør etableres konventioner for den kollaborative udvikling, så det er nemmere for andre at overtage opgaver ved behov.

\begin{itemize}
    \item \textbf{Uklare krav}
    \begin{itemize}
        \item[R1] Interessenter leverer utilstrækkelig eller modstridende feedback, hvilket vanskeliggør prioritering af opgaver i backlog'en.
        \item[R2] Begrænset repræsentation af interessenter medfører, at væsentlige krav først identificeres sent i projektforløbet.
    \end{itemize}

    \item \textbf{Teknisk gæld}
    \begin{itemize}
        \item[R3] Tidsmæssigt pres i sprint fører til, at refaktorering nedprioriteres i projektarbejdet.
        \item[R4] Uens teknisk niveau i projektgruppen medfører inkonsistent kodekvalitet og ekstra tid til code reviews.
    \end{itemize}


    \item \textbf{Projektorganisation}
    \begin{itemize}
        \item[R5] Et gruppemedlem bliver langtidssygemeldt, hvilket reducerer udviklingskapaciteten og forsinker projektet i.f.t. tidsplanen.
        \item[R6] Samtidig arbejdsbelastning fra andre fag reducerer den effektive arbejdstid i projektgruppen.
    \end{itemize}
\end{itemize}

\subsection{Risikomatrix}\label{subsec:risikomatrix}
Sandsynlighed og konsekvensgrad vurderes på en skala fra 1 til 3, hvor 1 er lav og 3 er høj.
En risikomatrix anvendes til at prioritere risici i udviklingsprocessen.
Farverne i risikomatrixen anvendes til at visualisere risikoniveauet, 
hvor grøn angiver lav risiko, gul angiver middel risiko, og rød angiver høj risiko.

\begin{table}[h]
    \centering
    \caption{5×5 risikomatrix}
    \label{tab:risikomatrix55}
    \begin{tabular}{|c|c|c|c|c|c|}
        \hline
         & \multicolumn{5}{|c|}{\textbf{Konsekvensgrad}} \\
        \hline
        \textbf{Sandsynlighed}
        & \textbf{Ubetydelig (1)}
        & \textbf{Lav (2)}
        & \textbf{Moderat (3)}
        & \textbf{Alvorlig (4)}
        & \textbf{Kritisk (5)} \\
        \hline
        \textbf{Meget lav (1)}
        & \cellcolor{green!40} 1
        & \cellcolor{green!40} 2
        & \cellcolor{green!40} 3
        & \cellcolor{yellow!40} 4
        & \cellcolor{yellow!40} 5 \\
        \hline
        \textbf{Lav (2)}
        & \cellcolor{green!40} 2
        & \cellcolor{green!40} 4
        & \cellcolor{yellow!40} 6
        & \cellcolor{yellow!40} 8
        & \cellcolor{red!40} 10 \\
        \hline
        \textbf{Middel (3)}
        & \cellcolor{green!40} 3
        & \cellcolor{yellow!40} 6
        & \cellcolor{yellow!40} 9 (R1)
        & \cellcolor{red!40} 12
        & \cellcolor{red!40} 15 \\
        \hline
        \textbf{Høj (4)}
        & \cellcolor{yellow!40} 4
        & \cellcolor{yellow!40} 8
        & \cellcolor{red!40} 12
        & \cellcolor{red!40} 16
        & \cellcolor{red!40} 20 \\
        \hline
        \textbf{Meget høj (5)}
        & \cellcolor{yellow!40} 5
        & \cellcolor{red!40} 10
        & \cellcolor{red!40} 15
        & \cellcolor{red!40} 20
        & \cellcolor{red!40} 25 \\
        \hline
    \end{tabular}
\end{table}


\section{Yderligere applikationsniveau risici}\label{sec:yderligere-risici}
Disse punkter er ikke relateret til risici i forbinde med udviklingsprocessen, men snarere risici forbundet med selve applikationen, dens succes og sælgbarhed som produkt.
Disse punkter indgår ikke i risikomatricen 

\subsubsection{Rolle- \& rettighedsstyring}
En anden væsentlig risiko er fejl i rolle- og rettighedsstyring, da platformen understøtter flere brugerroller med rettigheder der er både identitetsbaserede
og delvist overlappende.
Fejl her kan føre til manglende eller uautoriseret adgang til funktionaliteter.
Risikoen reduceres ved at implementere rollebaseret adgangskontrol (login) tidligt i udviklingen som en del af MVP og validere denne løbende gennem tests efter hver iteration.

\subsubsection{Datasikkerhed \& håndtering af persondata}
Datasikkerhed udgør ligeledes en risiko, da systemet potentielt behandler følsomme oplysninger om dets brugere.
Her er det afgørende, at der i processen laves en vurdering på baggrund af GDPR, og i den agile proces håndteres dette ved at indarbejde sikkerhed som en del af hvert sprint, så sikkerhedskrav implementeres inkrementelt frem for som en afsluttende aktivitet.

\subsubsection{Brugeroplevelse \& adoption}
Endelig er det meningen at applikationen skal levere en service, der overgår eksisterende løsninger i forhold til brugbarhed og lethed.
Brugerne har forskellig teknisk erfaring, hvor særligt gæster og almene medlemmer gerne skal være i stand til at navigere applikationen uden ekstern hjælp.
Hvis systemet opleves som komplekst kan det føre til lav anvendelse og ultimativt fravalg.
Denne risiko reduceres ved at stille eksplicitte krav til brugervenlighed, at føre løbende brugertests og få feedback fra interessenten i forbindelse med sprint reviews.
